\documentclass[paper=a4paper,fontsize=11pt]{jlreq}
\usepackage{amsmath,amssymb}

\title{リアルタイム\TeX プレビュー}
\author{TeXConf 2026 デモ}
\date{}

\begin{document}

\maketitle

\section{打った瞬間にページが変わる}

この文書は、常駐している Lua\LaTeX{} がその場で組んでいる。
左のソースを一文字直すと、その段落を含むブロックだけが組み直され、
変わったページだけが右側で差し替わる。行分割は本物の \TeX{} の
Knuth--Plass 法で決まり、和文の禁則処理や和欧文間の空白も
jlreq の設定どおりに入る。

段落の途中に数式 $e^{i\pi} + 1 = 0$ を書いても、見た目は本物の出力と
変わらない。数式の部分だけは、\TeX{} が組んだ本物の画像を上から重ねている。\footnote{脚注も本物と同じように、ページの下にまとめて置かれる。}

\section{ページを組むのは JavaScript}

ページ分割は、\TeX{} の規則をそのまま JavaScript で組み直している。
段落が一行増えれば、その行はきちんと次のページへ送られる。
\begin{equation}
  c = \begin{cases}
    p & (b < \infty,\ p \le -10000) \\
    b + p + q & (b < 10000,\ |p| < 10000) \\
    \infty & (b = \infty)
  \end{cases}
\end{equation}
上の式は、ページを区切る候補ごとのコストである。

\section{正しさは正本が保証する}

速い表示はあくまで仮のものである。裏では普通の lualatex が
文書全体を組んでいて、でき上がった正本の PDF と見比べ、
外れた所があれば描き方を一段慎重にする。

\begin{table}[htbp]
  \centering
  \begin{tabular}{lrr}
    \hline
    & 仮組み & フルコンパイル \\
    \hline
    1 文字打ってから & 約 35 ms & 約 1.3 s \\
    \hline
  \end{tabular}
  \caption{画面が変わるまでの時間}
\end{table}

\end{document}
